\section{Sufficient lifting for a fixed finite group}\label{sec:finite-lifting}

This section closes a specified version of the sufficient-lifting question left by K61. The source group $G$ is fixed and finite, the rank $d$ is fixed, the coefficient ring is $\mathbb Z_p$, and every lift must satisfy the same group multiplication table. We first exhibit a zero-exchange-defect representation that nevertheless does not lift. We then give a complete one-step obstruction and an explicit finite horizon after which convergent correction constructs an integral action. None of these assertions concerns arbitrary continuous representations of a changing system of finite Galois quotients.

\paragraph{Attribution.}
The one-step obstruction is the classical group-extension class in $H^2$, with the standard inhomogeneous conventions of \cite[Chapter II, 1.17--1.19]{MilneCFT2020}. Finite-group ultrametric stability and the quantitative bound used below are known: \cite[Lemma 7.3, Proposition 7.4, Example 7.5]{FournierFacio2026} supply the same valuation threshold, in a broader setting. We provide a direct matrix-averaging proof, translate it into an exact finite lifting criterion, and connect it to the present residual calculus. Historical novelty is not claimed for that bound or the deformation-theoretic mechanism.

\subsection{The relation residual and its repair class}

Put $R_n=\mathbb Z/p^n$, $k=\mathbb F_p$, and fix
\[
 \rho_n:G\longrightarrow\operatorname{GL}_d(R_n),\qquad n\geq1.
\]
Let $V=\operatorname{End}_k(k^d)$ with action
\[
 g\cdot X=\bar\rho(g)X\bar\rho(g)^{-1},
 \qquad \bar\rho=\rho_n\bmod p.
\]
Choose matrices $A_g\in\operatorname{GL}_d(R_{n+1})$ lifting $\rho_n(g)$, with $A_1=I$. They need not constitute a representation. Define the normalized relation residual
\[
 c_A(g,h)=\frac{A_gA_hA_{gh}^{-1}-I}{p^n}\pmod p\quad\in V.
\]
Division here identifies the square-zero ideal $p^nR_{n+1}$ with $k$. It is not inversion of $p$ in $R_{n+1}$. The square-zero assertion uses $2n\geq n+1$.

For a normalized $1$-cochain $b:G\to V$, set
\[
 (\delta b)(g,h)=b(g)+g\cdot b(h)-b(gh).
\]
For a normalized $2$-cochain set
\[
 (\delta c)(g,h,t)=g\cdot c(h,t)-c(gh,t)+c(g,ht)-c(g,h).
\]
A cochain is normalized when it vanishes whenever an argument is the identity.

\begin{theorem}[Complete one-step relation obstruction]\label{thm:relation-H2}
The residual $c_A$ is a $2$-cocycle. Its cohomology class
\[
 o_n(\rho_n)=[c_A]\in H^2(G,V)
\]
is independent of the chosen matrix lifts. Moreover,
\[
 \boxed{o_n(\rho_n)=0
 \quad\Longleftrightarrow\quad
 \rho_n\text{ lifts to }\operatorname{GL}_d(R_{n+1}).}
\]
More precisely, all lifts with the fixed reduction have the form
\[
 A_g^b=(I+p^n b(g))A_g,
\]
where representatives of $b(g)$ may be chosen arbitrarily, and the necessary and sufficient equation is
\[
 \boxed{\delta b=-c_A.}
\]
\end{theorem}
\begin{proof}
Associativity of the threefold product gives the exact identity
\[
 (I+p^nc_A(g,h))(I+p^nc_A(gh,t))
 = (I+p^n g\cdot c_A(h,t))(I+p^nc_A(g,ht))
\]
in $R_{n+1}$, where the conjugation can be reduced modulo $p$. All products of two $p^n$ terms vanish. Dividing the resulting linear identity by $p^n$ gives $\delta c_A=0$.

Any alternative matrix lift is uniquely of the form $(I+p^nb(g))A_g$, for a normalized cochain $b$. Expanding its multiplicativity residual and discarding square-zero terms gives
\[
 c_{A^b}=c_A+\delta b.
\]
Thus the class is unchanged and the new matrices satisfy every group relation exactly if and only if $\delta b=-c_A$. Such a cochain exists precisely when $[c_A]=0$. Invertibility is automatic because $I+p^nb(g)$ has inverse $I-p^nb(g)$ at this level.
\end{proof}

\begin{proposition}[Repair ambiguity, conjugacy, and finite certificates]\label{prop:repair-H1}
If the one-step lift set is nonempty, its framed lifts form a torsor for $Z^1(G,V)$. Conjugation by matrices $I+p^nX$, $X\in V$, quotients this torsor by coboundaries. Hence the strict-conjugacy classes form a torsor for $H^1(G,V)$, and the stabilizer of a lift under this conjugation is $H^0(G,V)$ as an additive group.

For a fixed multiplication table and residual action, deciding existence is a finite linear problem over $k$. A solution $b$ is a positive certificate. A linear functional $\lambda$ on normalized $2$-cochains with
\[
 \lambda\circ\delta=0,\qquad \lambda(c_A)\ne0
\]
is a negative certificate.
\end{proposition}
\begin{proof}
If $b_0$ is one solution, every other solution is $b_0+z$ with $\delta z=0$. Conjugating the matrices by $I+p^nX$ replaces $b(g)$ by $b(g)+X-g\cdot X$. These are exactly the negatives of the usual $0$-coboundaries and yield the same quotient subgroup. Conjugation fixes the action precisely when $X$ is invariant. The cochain spaces have dimensions $(|G|-1)d^2$ and $(|G|-1)^2d^2$. Gaussian elimination decides whether $-c_A$ is in the image of $\delta$. If it is not, elementary linear algebra supplies a functional annihilating that image but not $c_A$.
\end{proof}

This is an actual obstruction to realizing the finer group action, not merely failure of a selected invariant to commute with coefficient change. The repair changes matrices within their prescribed residue classes and keeps the coefficient unit. It does not take a quotient of the coefficient theory.

\subsection{A vanishing exchange cone can miss nonliftability}

\begin{theorem}[A central sign hides the obstruction from derived invariants]\label{thm:blind-exchange}
Let $p$ be odd, $R=\mathbb Z/p^2$, and let $J\in M_d(\mathbb Z)$ be invertible modulo $p$. On $M=R^d$ let
\[
 G=C_p\times C_2=\langle\sigma,\tau\rangle,
 \qquad \sigma\mapsto I+pJ,\quad \tau\mapsto-I.
\]
For the canonical comparison of \cref{thm:lifted-input-exchange},
\[
 \Gamma_R(M)\simeq0,
 \qquad \Gamma_k(M\otimes_R^{\mathbf L}k)\simeq0,
 \qquad \Cofib(\beta_{R,k,M})\simeq0.
\]
Nevertheless this representation does not lift to $\mathbb Z/p^3$, and hence does not lift as a finite free $\mathbb Z_p$-representation of this fixed group.
\end{theorem}
\begin{proof}
The two displayed matrices commute and satisfy the defining relations. Since $2$ is invertible in $R$, the central orthogonal idempotents
\[
 e_+=(1+\tau)/2,\qquad e_-=(1-\tau)/2
\]
split $R[G]$ into a product of rings. The trivial module $R$ is in the plus factor and $M$ in the minus factor. Every derived Hom between the two factors is zero, as is seen by resolving within either direct factor. The same decomposition persists after reduction to $k$. Since $M$ is free over $R$, its derived reduction is its ordinary reduction. Both endpoints of $\beta$ therefore vanish.

Any lift of this $G$-action would restrict to a lift of its $C_p$-action. By \cref{thm:one-step-group-lift}, each lift of $I+pJ$ has $p$th power $I+p^2J\ne I$ modulo $p^3$. This is impossible.
\end{proof}

The example refutes sufficiency of the single $\Gamma$-comparison test in the finite-group class. It does not merely leave the converse unproved. The loss is a loss of detection: a functor can see none of an input while the input still has inconsistent lifting relations. In contrast, conjugation by $-I$ is trivial on $\operatorname{End}_k(k^d)$, so the adjoint relation obstruction is not hidden by this sign. Twisting by the integral sign character also restores the earlier nonzero exchange example. Neither observation asserts that a finite family of such probe comparisons classifies all lifting obstructions.

\subsection{Quadratic correction of an approximate action}

Let $N=|G|$ and write $e=v_p(N)$. For a matrix $X$ over $\mathbb Q_p$, the assertion $v_p(X)\geq a$ means every entry lies in $p^a\mathbb Z_p$. Invertible integral matrices and their inverses preserve this divisibility under multiplication. Let $A:G\to\operatorname{GL}_d(\mathbb Z_p)$ be a normalized map, not necessarily a homomorphism. Define
\[
 E(g,h)=A_gA_hA_{gh}^{-1}-I,
 \qquad \alpha_g(X)=A_gXA_g^{-1}.
\]
Suppose $E(g,h)\in p^mM_d(\mathbb Z_p)$ for every pair. The operators $\alpha_g$ are not assumed to give a group action. This distinction is essential before the repair has succeeded.

\begin{lemma}[Averaged correction with controlled precision loss]\label{lem:quadratic-repair}
If $m>e$, define in $M_d(\mathbb Q_p)$
\[
 b(g)=\frac1N\sum_{t\in G}E(g,t),\qquad
 A'_g=(I-b(g))A_g.
\]
Then $b(g)\in p^{m-e}M_d(\mathbb Z_p)$, $b(1)=0$, and $A'$ is a normalized map into $\operatorname{GL}_d(\mathbb Z_p)$. Its defect satisfies
\[
 \boxed{A'_gA'_h(A'_{gh})^{-1}-I\in
 p^{2m-2e}M_d(\mathbb Z_p).}
\]
Also $A'_g-A_g\in p^{m-e}M_d(\mathbb Z_p)$.
\end{lemma}
\begin{proof}
Division by $N=p^eu$, $u\in\mathbb Z_p^\times$, loses at most $e$ digits, proving the bound on $b$. Normalization gives $b(1)=0$. Because $m-e\geq1$, each $I-b(g)$ is invertible.

Associativity gives an identity without assuming an exact adjoint action:
\[
 (I+E(g,h))(I+E(gh,t))
 = (I+\alpha_g E(h,t))(I+E(g,ht)).
\]
Consequently
\[
 E(g,h)+E(gh,t)-\alpha_g E(h,t)-E(g,ht)
 \in p^{2m}M_d(\mathbb Z_p).
\]
Sum over $t$ and divide by $N$. The permutation $t\mapsto ht$ identifies the last average with $b(g)$, giving
\[
 E(g,h)-b(g)-\alpha_g b(h)+b(gh)
 \in p^{2m-e}M_d(\mathbb Z_p).
\]
Put $a=m-e>0$. The repaired multiplicativity product is exactly
\[
 (I-b(g))(I-\alpha_g b(h))(I+E(g,h))(I-b(gh))^{-1}.
\]
Modulo $p^{2a}$ its difference from $I$ is the linear expression just displayed: products of two corrections vanish at that precision, and $(I-b)^{-1}\equiv I+b\pmod{p^{2a}}$. That expression is divisible by $p^{2m-e}$, which is contained in $p^{2a}$. This proves the new defect bound. The change in $A_g$ is $-b(g)A_g$ and has the asserted precision.
\end{proof}

This proof exhibits the residual as input to a repair operator. The correction is admissible because its denominators are absorbed by the defect's valuation. It requires no prior integral action and does not delete any module or coefficient information.

\begin{theorem}[Convergent finite-group action repair]\label{thm:convergent-action-repair}
If the initial defect depth satisfies $m>2e$, the preceding averaging operation iterates to a representation
\[
 \rho:G\longrightarrow\operatorname{GL}_d(\mathbb Z_p)
 \quad\text{with}\quad
 \rho(g)\equiv A_g\pmod{p^{m-e}}.
\]
After $j$ iterations the certified defect depth is
\[
 m_j=2e+2^j(m-2e),
\]
and the next change has depth at least $m_j-e$. In particular, the sequence converges uniformly on $G$ and all prescribed multiplication relations hold exactly in the limit.
\end{theorem}
\begin{proof}
Set $m_0=m$ and $m_{j+1}=2m_j-2e$. Induction gives the displayed formula, so $m_j>2e$ at every stage and $m_j\to\infty$. The correction lemma applies repeatedly. Differences of consecutive matrices have entry valuations tending to infinity, hence their telescoping sums converge in the complete ring $M_d(\mathbb Z_p)$. Every correction preserves reduction modulo $p$, so the limits remain invertible. Their reduction modulo $p^{m-e}$ is the original one.

Matrix multiplication and inversion on $\operatorname{GL}_d(\mathbb Z_p)$ are continuous. The defects converge to zero because their depths are at least $m_j$. Passing to the limit gives $\rho(g)\rho(h)=\rho(gh)$ and $\rho(1)=I$.
\end{proof}

The recurrence is a quantitative certificate of convergence, not a claim that arbitrary small residuals can always be removed. At $m=2e$ this estimate does not improve the defect depth. We do not assert that the threshold is optimal for every group, rank, or representation.

\subsection{A finite necessary-and-sufficient lifting horizon}

\begin{theorem}[Finite horizon for fixed finite-group representations]\label{thm:finite-action-horizon}
Let $G$ be a finite group, $d,n\geq1$, and $e=v_p(|G|)$. Set
\[
 \boxed{h_G(n)=\max\{n+e,\,2e+1\}.}
\]
For every $\rho_n:G\to\operatorname{GL}_d(\mathbb Z/p^n)$, the following are equivalent:
\begin{enumerate}
\item There is a representation $\rho:G\to\operatorname{GL}_d(\mathbb Z_p)$ with reduction $\rho_n$.
\item There is a representation $\rho_h:G\to\operatorname{GL}_d(\mathbb Z/p^{h_G(n)})$ with reduction $\rho_n$.
\end{enumerate}
The implication from (2) to (1) is constructive to arbitrary requested precision by \cref{thm:convergent-action-repair}. It preserves the original level $n$, not necessarily every digit of the chosen level-$h$ witness.
\end{theorem}
\begin{proof}
Reduction proves (1) implies (2). Conversely, lift every matrix of $\rho_h$ arbitrarily to an invertible integral matrix, keeping the identity equal to $I$. Its multiplicativity defects are divisible by $p^h$. Since $h>2e$, convergent repair applies and returns an action agreeing with these matrices modulo $p^{h-e}$. Because $h-e\geq n$, it returns precisely the specified $\rho_n$ on reduction.
\end{proof}

Thus a supplied high-precision action is a finite certificate for existence of an entire integral action. A level-$h$ witness may itself have an obstruction at its next level: the repair is allowed to alter the top $e$ digits while keeping the requested lower data. There is no claim that a positive certificate yields an exact lift of that witness at every digit.

A terminating decision method follows. Begin with $\rho_n$ and use \cref{thm:relation-H2,prop:repair-H1} to enumerate all one-step solutions, retaining their actual framed matrices, until level $h_G(n)$. A surviving branch is a positive witness. If all branches die before that finite level, no integral lift exists. Every layer is finite, and Gaussian elimination provides checkable obstruction certificates at dead nodes. Choosing a single convenient branch is not sufficient for a negative verdict. Quotienting branches by conjugacy is optional and requires preserving their lifting data.

For $p\nmid|G|$, one has $e=0$ and $h_G(n)=n$. Every finite-level representation in this fixed finite-group problem then lifts. When $e>0$, the test is still finite, but its number of branches may be enormous. No polynomial-time or uniformly practical algorithm is claimed.

\begin{example}[A positive witness need not lift itself]\label{ex:horizon-witness}
Take $p=2$, $G=C_2$, and rank one. The level-two representation with generator $3\pmod4$ has the level-three witness $3\pmod8$, since $3^2\equiv1\pmod8$. Here $e=1$ and $h_G(2)=3$. Nevertheless neither lift of $3\pmod8$ to $\mathbb Z/16$ has square one: both $3^2$ and $11^2$ are $9\pmod{16}$.

Starting with the integer matrix $a_0=3$, the averaged repair becomes
\[
 a_{j+1}=\frac{3-a_j^2}{2}\,a_j.
\]
The first values are $3,-9,351$, whose relation defects have valuations $3,4,6$. The convergence theorem gives an exact $2$-adic involution congruent to $3\pmod4$. In the integral domain $\mathbb Z_2$, the equation $a^2=1$ has only the solutions $1,-1$. Thus the repaired limit is $-1$, which is $7\pmod8$, not the selected witness $3\pmod8$. This is exactly the distinction between preserving the requested seed and preserving every exploratory refinement.
\end{example}

\subsection{Controlled uniqueness and a finite catalogue}

\begin{theorem}[Close integral actions are conjugate]\label{thm:close-action-conjugacy}
Let $\rho,\sigma:G\to\operatorname{GL}_d(\mathbb Z_p)$ be representations satisfying
\[
 \sigma(g)\equiv\rho(g)\pmod{p^s}\quad(g\in G),\qquad s>e.
\]
Then
\[
 T=\frac1{|G|}\sum_{g\in G}\sigma(g)\rho(g)^{-1}
 \in I+p^{s-e}M_d(\mathbb Z_p)
\]
is invertible and satisfies
\[
 \sigma(g)T=T\rho(g),\qquad
 \sigma(g)=T\rho(g)T^{-1}.
\]
\end{theorem}
\begin{proof}
Every summand is congruent to $I$ modulo $p^s$. Its average is therefore $I$ modulo $p^{s-e}$, so it is integral and invertible. For $h\in G$,
\[
 \sigma(h)T\rho(h)^{-1}
 =\frac1{|G|}\sum_g\sigma(hg)\rho(hg)^{-1}=T,
\]
by reindexing the finite sum. This is the required intertwining identity.
\end{proof}

This is the averaging mechanism underlying finite-precision rigidity of integral finite-group representations, closely related to the classical rigidity of integral finite-group lattices. The proof above is independent of any such external rigidity theorem. The strict inequality matters: for $p=2$, $G=C_2$, and rank one, the trivial and sign representations agree modulo $2$ but are not conjugate. Here $s=e=1$.

\begin{corollary}[A finite catalogue of integral action classes]\label{cor:finite-action-catalogue}
Let $s=e+1$ and $h=2e+1$. Denote by $\mathscr L_s$ the image of
\[
 \operatorname{Hom}(G,\operatorname{GL}_d(\mathbb Z/p^h))
 \longrightarrow
 \operatorname{Hom}(G,\operatorname{GL}_d(\mathbb Z/p^s)).
\]
Reduction gives a bijection
\[
 \boxed{
 \frac{\operatorname{Hom}(G,\operatorname{GL}_d(\mathbb Z_p))}
 {\operatorname{GL}_d(\mathbb Z_p)}
 \ \longleftrightarrow\
 \frac{\mathscr L_s}{\operatorname{GL}_d(\mathbb Z/p^s)}.}
\]
Consequently this class set is finite and can be determined from a finite multiplication table and finite matrix data.
\end{corollary}
\begin{proof}
The finite-horizon theorem gives exactly $h_G(s)=2e+1$, so every element of $\mathscr L_s$ lifts and every integral action reduces into it. If two reductions are conjugate, lift their conjugating matrix to $\operatorname{GL}_d(\mathbb Z_p)$ and first conjugate one action by that lift. Their reductions now agree modulo $p^s$. Since $s>e$, \cref{thm:close-action-conjugacy} makes the integral actions conjugate. The target is finite because it is a quotient of a finite set.
\end{proof}

This is a classification of isomorphism classes, not an equivalence with a finite groupoid. Integral automorphism groups and noninvertible morphisms need not be determined by finite reduction. Likewise the averaging construction is not asserted to be a canonical framed lift. Two outputs agreeing modulo $p^{m-e}$ are conjugate by a matrix congruent to $I$ modulo $p^{m-2e}$ when $m>2e$, but need not be literally the same matrices.

\subsection{A smaller certificate for cyclic groups}

\begin{proposition}[Cyclic relation Jacobian]\label{prop:cyclic-jacobian}
Let $G=C_q$ for an integer $q\geq1$, and let $A\in\operatorname{GL}_d(\mathbb Z/p^n)$ satisfy $A^q=I$. Choose a lift $\widetilde A$ modulo $p^{n+1}$ and set
\[
 c=\frac{\widetilde A^q-I}{p^n}\pmod p,
 \qquad L_{\bar A}(B)=\sum_{i=0}^{q-1}\bar A^iB\bar A^{q-1-i}.
\]
Then the lifts of the cyclic action to level $n+1$ are exactly the matrices
\[
 \widetilde A+p^nB\quad\text{such that}\quad L_{\bar A}(B)=-c.
\]
If the affine solution set is nonempty, it has $p^{\dim_k\ker L_{\bar A}}$ elements. This criterion uses $d^2$ variables instead of all group-table cochains.
\end{proposition}
\begin{proof}
Expand the noncommutative product $(\widetilde A+p^nB)^q$. Terms with two or more factors $p^nB$ vanish modulo $p^{n+1}$. The terms with one such factor give exactly $p^nL_{\bar A}(B)$. Therefore the cyclic relation is equivalent to the stated linear equation. All candidate matrices remain invertible because their mod-$p$ reductions are fixed. A nonempty fiber of a linear map over $k$ is a translate of its kernel and has the claimed cardinality.
\end{proof}

For $A=I+pJ$ modulo $p^2$ and $q=p$, the residual is $c=J$ while $L_I(B)=pB=0$ in $k$. This recovers the one-step obstruction without exhaustive testing. At other inputs the Jacobian may have nonzero rank, so some nonzero initial relation residuals are removable rather than persistent.

\subsection{Why a fixed finite quotient is a real restriction}

\begin{proposition}[A profinite lift need not preserve the finite quotient]\label{prop:profinite-quotient-boundary}
Let $p$ be odd and $J\in M_d(\mathbb Z_p)$ be invertible modulo $p$. The action $C_p\to\operatorname{GL}_d(\mathbb Z/p^2)$ given by $I+pJ$ has no integral lift as an action of $C_p$. Its pullback along $\mathbb Z_p\twoheadrightarrow C_p$, however, has a continuous integral lift
\[
 \rho(x)=(I+pJ)^x
 =\sum_{j\geq0}\binom{x}{j}p^jJ^j,
 \qquad x\in\mathbb Z_p.
\]
\end{proposition}
\begin{proof}
For each $j$, the polynomial $\binom{x}{j}$ is continuous on $\mathbb Z_p$ and takes values there: approximate $x$ by nonnegative integers, where the binomial coefficient is integral, and use closedness of $\mathbb Z_p$ in $\mathbb Q_p$. Thus the series converges uniformly because its $j$th term is divisible by $p^j$. For nonnegative integer arguments it is the usual matrix power. The nonnegative integers are dense in $\mathbb Z_p$, so continuity extends the identity $\rho(x+y)=\rho(x)\rho(y)$ to all pairs. Each value is $I$ modulo $p$ and hence invertible.

Modulo $p^2$ the series is $I+pxJ$, which depends only on $x\bmod p$ and gives the specified finite-level action. But $\rho(p)\equiv I+p^2J\ne I\pmod{p^3}$, so the integral lift does not factor through $C_p$. Nonexistence of an integral $C_p$-lift is \cref{thm:one-step-group-lift}.
\end{proof}

The same warning applies to finite quotients used in a Galois problem. A relation imposed by fixing one finite quotient may be absent from the full profinite source. Failure to lift through that fixed quotient therefore does not by itself obstruct a continuous lift with growing image. Conversely, the finite-group horizon depends on $v_p(|G|)$ and gives no uniform bound for arbitrary growing finite quotients.

\subsection{What the new block establishes}

There are now three distinct tests. A chosen functor's exchange cone is a necessary diagnostic on its detectable inputs, but the central-sign example shows it is not a complete lifting invariant. The class $o_n(\rho_n)$ is a complete one-step obstruction because its vanishing is proved equivalent to solving the actual multiplicativity equations. Finally, the finite horizon plus convergent correction turns a finite witness into a full $\mathbb Z_p$-action for a fixed finite group. This gives a positive reconstruction theorem without killing the residual in a Verdier quotient.

In the two-axis vocabulary, the resolution fiber is the actual lift groupoid, the transformation direction includes changes of basis, the relation residual is a normalized cocycle, admissible repairs are cochains, and its class in $H^2$ is the obstruction. The averaging theorem then supplies a quantitative regime in which repeated repairs converge. These identifications recover and connect classical deformation and ultrametric-stability mathematics. They do not make the abstract obstruction formalism a proof of automorphic existence or Langlands reciprocity.
